% =============================================================================
% Section: Seven Non-Negotiable Doctrines
% =============================================================================

\section{Seven Non-Negotiable Doctrines}
\label{sec:doctrines}

The three primitives ($\Del$, $\Pit$, $\Tim$) and the derivation from
nothingness impose seven constraints on any system that reasons under A0$^*$.
These are not design choices or engineering heuristics---they are forced
consequences of the axiom.  Each doctrine is traced to the derivation step
that forces it.

\begin{enumerate}[label=\textbf{D\arabic*.}]
  \item \textbf{Everything must be finitely describable.}
        Every admissible object is a finite binary string in $\Carrier$.
        Infinite descriptions require infinite witnesses, violating A0$^*$.
        \hfill\emph{(A0$^*$, Step~1; \cref{thm:carrier})}

  \item \textbf{Meaning requires testable differences.}
        A distinction without a separating test $\tau \in \Del^*$ is
        operationally void.  If no finite procedure can detect the
        difference, the difference does not exist within the framework.
        \hfill\emph{(Steps~3--4; \cref{thm:executability}, \cref{thm:endogenous})}

  \item \textbf{Truth emerges from feasible tests.}
        The truth quotient $\TQ{\Ledger}$ is the coarsest structure
        surviving all executed tests.  Truth is not assumed; it is
        constructed by the accumulation of test outcomes.
        \hfill\emph{(Steps~6--7; \cref{thm:truth-quotient}, \cref{thm:indistinguishability})}

  \item \textbf{Time is an irreversible ledger.}
        Each record appended to $\Ledger$ advances the clock.  Reversal
        would destroy a witness, violating A0$^*$.  Time is the irreversibility
        of record formation.
        \hfill\emph{(Step~8; \cref{thm:time})}

  \item \textbf{The truthpoint $\Truthpoint$ exists as a limit concept.}
        The ideal truth quotient under unlimited budget is approached but
        never operationally reached.  The gap between any finite
        $\TQ{\Ledger}$ and $\Truthpoint$ is precisely what $\OMEGA$ reports.
        \hfill\emph{(Step~11; \cref{thm:opnoth})}

  \item \textbf{Completeness demands enumeration of lawful tests.}
        The canonical enumeration $\CanEnum$ over $\Del^*$ ensures no
        admissible test is overlooked.  The enumeration order is determined
        by the self-delimiting syntax $\SdMap$, requiring no external
        scheduler.
        \hfill\emph{(Step~4; \cref{thm:endogenous})}

  \item \textbf{Claims require finite witnesses or explicit gaps.}
        Every output is either $\UNIQUE$ (verified answer with witness) or
        $\OMEGA$ (unknown, with the minimal separator test needed for
        resolution).  There is no probabilistic hedging, no confidence
        score, no ungrounded assertion.
        \hfill\emph{(Step~0; \cref{thm:nothingness})}
\end{enumerate}

\begin{remark}
These doctrines are \emph{consequences} of A0$^*$ and the derivation, collected
here for clarity.  They add no axiomatic content beyond what the derivation
establishes.  Any system that violates a doctrine necessarily violates A0$^*$ at
the cited step.
\end{remark}
